\newif\ifglobal

%%%%%%%%%%%%%%%%%%%%%%%
%%% Compile to view %%%
%%%%%%%%%%%%%%%%%%%%%%%

%%%%%%%%%%%%%%%%%%%%%%%%%%%%%%%%%%%%%%%%%%%%%%%%%%%
%%% Readme:                                     %%%
%%% https://www.overleaf.com/read/bwdjvfjksvjy/ %%%
%%%%%%%%%%%%%%%%%%%%%%%%%%%%%%%%%%%%%%%%%%%%%%%%%%%

% >>> M?: marks a module setting number or important notice

% >>> M4: Set {./relative-path} to external files for this module <<<
\def \modulefiles {./27-SHA}
% To add an external file, use:
% (Replace '---' with actual filename)
% '\input{\modulefiles/tables/---.tex}' for tables
% '\includegraphics{\modulefiles/figures/---}' for figures
% '\subfile{\modulefiles/---}' for register files

% >>> M1: This module is in global project? <<<
% 'YES' - uncomment; 'NO' - comment out
\globaltrue

\ifglobal
    \documentclass[main\_\_EN.tex]{subfiles}

\else
    \documentclass[openany, 10pt]{book}

    % >>> M2: In ./00-shared/preamble-trm-module, also find
    %         settings for visibility of todo notes, watermark <<<
    \usepackage{preamble-trm-module}

    % >>> M3: Temporary labels in English,
    %         you can add your own labels there <<<
    \myexternaldocument{./00-shared/temp-labels-en}

\fi %\ifglobal


\setcounter{secnumdepth}{4}

% >>> M5: If this module needs any updates in
% './00-shared/preamble-trm-module.sty'
% (include additional package, etc.), contact your module PM
% or create the file './00-shared/changelog.tex' make the
% necessary updates yourself and document them in this file <<<


\begin{document}


% Sets language into which pre-defined bits of text are translated
\selectlanguage{english}

% Do not delete this block
\ifglobal
% Add the following front matter if
% this module is outside of global project
\else
    \InputIfFileExists{readme}{}{}
    {\let\clearpage\relax \listoftodos}
    \tableofcontents
    %\listoftables
    %\listoffigures
\fi


%%%%%%%%%%%%%%%%%%%%%%%%%
%%% TRM Module Begins %%%
%%%%%%%%%%%%%%%%%%%%%%%%%

\hypertarget{sha}{}
\chapter{SHA Accelerator (SHA)}
\label{mod:sha}
\section{Introduction}

\chipname{} integrates an SHA accelerator, which is a hardware device that speeds up SHA algorithm significantly, compared to SHA algorithm implemented solely in software. The SHA accelerator integrated in \chipname{} has two working modes, which are  \hyperref[sec:sha-typical-sha]{Typical SHA} and  \hyperref[sec:sha-dma-sha]{DMA-SHA}.

\section{Features}

The following functionality is supported:

\begin{itemize}
    \item All the hash algorithms introduced in   \href{https://doi.org/10.6028/NIST.FIPS.180-4}{FIPS PUB 180-4 Spec}.
    \begin{itemize}
        \item SHA-1
        \item SHA-224
        \item SHA-256
        \item SHA-384
        \item SHA-512
        \item SHA-512/224
        \item SHA-512/256
        \item SHA-512/{\it{t}}
    \end{itemize}
    \item Two working modes
    \begin{itemize}
        \item Typical SHA
        \item DMA-SHA
    \end{itemize}
    \item interleaved function when working in Typical SHA working mode
    \item Interrupt function when working in DMA-SHA working mode
\end{itemize}

\section{Working Modes}
The SHA accelerator integrated in \chipname{} has two working modes.

\begin{itemize}
    \item \hyperref[sec:sha-typical-sha]{Typical SHA Working Mode}: all the data is written and read via CPU directly.
    \item \hyperref[sec:sha-dma-sha]{DMA-SHA Working Mode}: all the data is read via DMA. That is, users can configure the DMA controller to read all the data needed for hash operation, thus releasing CPU for completing other tasks.
\end{itemize}

Users can start the SHA accelerator with different working modes by configuring registers \hyperref[regdesc:SHASTARTREG]{SHA\_START\_REG} and \hyperref[regdesc:SHADMASTARTREG]{SHA\_DMA\_START\_REG}. For details, please see Table \ref{tab:sha-working-mode}.

\begin{longtable}[c]{ | l | l | }
\caption{SHA Accelerator Working Mode}
\label{tab:sha-working-mode}
\\\hline
\rowcolor{lightgray}
Working Mode      & Configuration Method\\
    \hline
    \hyperref[sec:sha-typical-sha]{Typical SHA}      &  Set \hyperref[regdesc:SHASTARTREG]{SHA\_START\_REG} to 1   \\
    \hline
    \hyperref[sec:sha-dma-sha]{DMA-SHA}           &  Set \hyperref[regdesc:SHADMASTARTREG]{SHA\_DMA\_START\_REG} to 1   \\
    \hline
\end{longtable}

Users can choose hash algorithms by configuring the \hyperref[regdesc:SHAMODEREG]{SHA\_MODE\_REG} register. For details, please see Table \ref{tab:sha-hash-algo-select}.

\begin{longtable}[c]{| l | c |}
\caption{SHA Hash Algorithm Selection}
\label{tab:sha-hash-algo-select}
\\\hline
\rowcolor{lightgray}
Hash Algorithm & \hyperref[regdesc:SHAMODEREG]{SHA\_MODE\_REG} Configuration\\
    \hline
    SHA-1 & 0\\
    \hline
    SHA-224 & 1\\
    \hline
    SHA-256 & 2\\
    \hline
    SHA-384 & 3\\
    \hline
    SHA-512 & 4\\
    \hline
    SHA-512/224 & 5\\
    \hline
    SHA-512/256 & 6\\
    \hline
    SHA-512/{\it{t}} & 7\\
    \hline
\end{longtable}

\begin{importantlisting}
\vspace{-0.5em}

    \chipname{}'s \nameref{mod:digsig} and \nameref{mod:hmac} modules also call the SHA accelerator. Therefore, users cannot access the SHA accelerator when these modules are working.
\end{importantlisting}

\section{Function Description}

SHA accelerator can generate the message digest via two steps: \hyperref[sec:sha-preprocessing]{Preprocessing} and  \hyperref[sec:sha-hash-operation]{Hash operation}.

\subsection{Preprocessing}\label{sec:sha-preprocessing}

Preprocessing consists of three steps: \hyperref[sec:sha-padding-msg]{padding the message}, \hyperref[sec:sha-parsing-msg]{parsing the message into message blocks} and \hyperref[sec:sha-set-initial-hash]{setting the initial hash value}.

\subsubsection{Padding the Message}\label{sec:sha-padding-msg}

The SHA accelerator can only process message blocks of 512 or 1024 bits, depending on the algorithm. Thus, all the messages should be padded to a multiple of 512 or 1024 bits before the hash task.

Suppose that the length of the message \begin{math}M\end{math} is \begin{math} m \end{math} bits. Then \begin{math} M \end{math} shall be padded as introduced below:

\begin{itemize}
    \item \textbf{SHA-1}, \textbf{SHA-224} and \textbf{SHA-256}
        \begin{enumerate}
        \item First, append the bit ``1" to the end of the message;
        \item Second, append \begin{math} k \end{math} zero bits, where \begin{math} k \end{math} is the smallest, non-negative solution to the equation \begin{math}m + 1 + k \ {\equiv} \ 448 \ mod \end{math} 512;
        \item Last, append the 64-bit block of value equal to the number \begin{math} m \end{math} expressed using a binary representation.
        \end{enumerate}
    \item \textbf{SHA-384}, \textbf{SHA-512}, \textbf{SHA-512/224}, \textbf{SHA-512/256} and \textbf{SHA-512/{\it{t}}}
        \begin{enumerate}
        \item First, append the bit ``1" to the end of the message;
        \item Second, append \begin{math} k \end{math} zero bits, where \begin{math} k \end{math} is the smallest, non-negative solution to the equation \begin{math}m + 1 + k \ {\equiv} \ 896 \ mod \end{math} 1024;
        \item Last, append the 128-bit block of value equal to the number \begin{math} m \end{math} expressed using a binary representation.
        \end{enumerate}
\end{itemize}

For more details, please refer to Section ``5.1 Padding the Message" in \href{https://doi.org/10.6028/NIST.FIPS.180-4}{FIPS PUB 180-4 Spec}.

\subsubsection{Parsing the Message}\label{sec:sha-parsing-msg}

The message and its padding must be parsed into \begin{math}N\end{math} 512-bit or 1024-bit blocks.

\begin{itemize}
    \item For \textbf{SHA-1}, \textbf{SHA-224} and \textbf{SHA-256}: the message and its padding are parsed into \begin{math}N\end{math} 512-bit blocks, \begin{math}M^{(1)}\end{math}, \begin{math}M^{(2)}\end{math}, \dots, \begin{math}M^{(N)}\end{math}. Since the 512 bits of the input block may be expressed as sixteen 32-bit words, the first 32 bits of message block \begin{math}i\end{math} are denoted M$^{(i)}_{0}$, the next 32 bits are M$^{(i)}_{1}$, and so on up to M$^{(i)}_{15}$.
    \item For \textbf{SHA-384}, \textbf{SHA-512}, \textbf{SHA-512/224}, \textbf{SHA-512/256} and \textbf{SHA-512/{\it{t}}}: the message and its padding are parsed into \begin{math}N\end{math} 1024-bit blocks. Since the 1024 bits of the input block may be expressed as sixteen 64-bit words, the first 64 bits of message block \begin{math}i\end{math} are denoted M$^{(i)}_{0}$, the next 64 bits are M$^{(i)}_{1}$, and so on up to M$^{(i)}_{15}$.
\end{itemize}

During the task, all the message blocks are written into the \hyperref[regdesc:SHAMNREG]{SHA\_M\_\regindex{n}\_REG}, following the rules below:
\begin{itemize}
    \item For \textbf{SHA-1, SHA-224 and SHA-256}: M$^{(i)}_{0}$ is stored in \hyperref[regdesc:SHAMNREG]{SHA\_M\_0\_REG}, M$^{(i)}_1$ stored in \hyperref[regdesc:SHAMNREG]{SHA\_M\_1\_REG}, \dots, and M$^{(i)}_{15}$ stored in \hyperref[regdesc:SHAMNREG]{SHA\_M\_15\_REG}.
    \item For \textbf{SHA-384, SHA-512, SHA-512/224 and SHA-512/256}: the most significant 32 bits and the least significant 32 bits of M$^{(i)}_0 $ are stored in \hyperref[regdesc:SHAMNREG]{SHA\_M\_0\_REG} and \hyperref[regdesc:SHAMNREG]{SHA\_M\_1\_REG}, respectively, \dots, the most significant 32 bits and the least significant 32 bits of M$^{(i)}_{15}$ are stored in \hyperref[regdesc:SHAMNREG]{SHA\_M\_30\_REG} and \hyperref[regdesc:SHAMNREG]{SHA\_M\_31\_REG}, respectively.

\end{itemize}

\begin{tiplisting}
\vspace{-0.5em}
For more information about ``message block", please refer to Section ``2.1 Glossary of Terms and Acronyms" in \href{https://doi.org/10.6028/NIST.FIPS.180-4}{FIPS PUB 180-4 Spec}.
\end{tiplisting}

\subsubsection{Initial Hash Value}\label{sec:sha-set-initial-hash}
Before hash task begins for each of the secure hash algorithms, the initial Hash value H$^{(0)}$ must be set based on different algorithms, among which the SHA-1, SHA-224, SHA-256, SHA-384, SHA-512, SHA-512/224, and SHA-512/256 algorithms use the initial Hash values (constant C) stored in the hardware.

However, SHA-512/{\it{t}} requires a distinct initial hash value for each operation for a given value of \begin{math}t\end{math}. Simply put, SHA-512/{\it{t}} is the generic name for a \begin{math}t\end{math}-bit hash function based on SHA-512 whose output is truncated to \begin{math}t\end{math} bits. \begin{math}t\end{math} is any positive integer without a leading zero such that \begin{math}t\end{math}<512, and \begin{math}t\end{math} is not 384. The initial hash value for SHA-512/{\it{t}} for a given value of {\it{t}} can be calculated by performing SHA-512 from hexadecimal representation of the string ``SHA-512/{\it{t}}". It's not hard to observe that when determining the initial hash values for SHA-512/{\it{t}} algorithms with different {\it{t}}, the only difference lies in the value of {\it{t}}.

Therefore, we have specially developed the following simplified method to calculate the initial hash value for SHA-512/{\it{t}}:

\begin{enumerate}\label{others:sha-512/t}
    \item \textbf{Generate t\_string and t\_length}: t\_string is a 32-bit data that stores the input message of {\it{t}}. t\_length is a 7-bit data that stores the length of the input message. The t\_string and t\_length are generated in methods described below, depending on the value of \begin{math}t\end{math}:

    \begin{itemize}
        \item If $1<=t<=9$, then $t\_length=7'h48$ and $t\_string$ is padded in the following format:\\
        \begin{longtable}[c]{|l|l|l|}
        \hline
        $8'h30+8'ht_0$ & $1'b1$ & $23'b0$\\
        \hline
        \end{longtable}
        where $t_0=t$.

        For example, if {\it{t}} = 8, then $t_0$ = 8 and $t\_string$ = $32'h38800000$.

        \item If $10<=t<=99$, then $t\_length=7'h50$ and $t\_string$ is padded in the following format:\\

        \begin{longtable}[c]{|l|l|l|l|}
        \hline
        $8'h30+8'ht_1$ & $8'h30+8'ht_0$ & $1'b1$ & $15'b0$\\
        \hline
        \end{longtable}
        where, $t_0=t\%10$ and $t_1=t/10$.

        For example, if {\it{t}} = 56, then $t_0$ = 6, $t_1$ = 5, and $t\_string$ = $32'h35368000$.

        \item If $100<=t<512$, then $t\_length=7'h58$ and $t\_string$ is padded in the following format:\\

        \begin{longtable}[c]{|l|l|l|l|l|}
        \hline
        $8'h30+8'ht_2$ & $8'h30+8'ht_1$ & $8'h30+8'ht_0$ & $1'b1$ & $7'b0$\\
        \hline
        \end{longtable}
        where, $t_0=t\%10$, $t_1=(t/10)\%10$, and $t_2=t/100$.

        For example, if {\it{t}} = 231, then $t_0$ = 1, $t_1$ = 3, $t_2$ = 2, and $t\_string$ = $32'h32333180$.

    \end{itemize}
    \item \textbf{Initialize relevant registers}: Initialize \hyperref[regdesc:SHATSTRINGREG]{SHA\_T\_STRING\_REG} and  \hyperref[regdesc:SHATLENGTHREG]{SHA\_T\_LENGTH\_REG} with the generated t\_string and t\_length in the previous step.

    \item \textbf{Obtain initial hash value}: Set the \hyperref[regdesc:SHAMODEREG]{SHA\_MODE\_REG} register to 7. Set the \hyperref[regdesc:SHASTARTREG]{SHA\_START\_REG} register to 1 to start the SHA accelerator. Then poll register \hyperref[regdesc:SHABUSYREG]{SHA\_BUSY\_REG} until the content of this register becomes 0, indicating the calculation of initial hash value is completed.
\end{enumerate}

Please note that the initial value for SHA-512/{\it{t}} can be also calculated according to the Section ``5.3.6 SHA-512/{\it{t}}" in \href{https://doi.org/10.6028/NIST.FIPS.180-4}{FIPS PUB 180-4 Spec}, that is performing SHA-512 operation (with its initial hash value set to the result of 8-bitwise XOR operation of C and 0x{}a5) from the hexadecimal representation of the string ``SHA-512/{\it{t}}".

\subsection{Hash Task Process}\label{sec:sha-hash-operation}

After the preprocessing, the \chipname{} SHA accelerator starts to hash a message \begin{math}M\end{math} and generates message digest of different lengths, depending on different hash algorithms. As described above, the \chipname{} SHA accelerator supports two working modes, which are \hyperref[sec:sha-typical-sha]{Typical SHA} and \hyperref[sec:sha-dma-sha]{DMA-SHA}. The
operation process for the SHA accelerator under two working modes is described in the following subsections.

\subsubsection{Typical SHA Mode Process}\label{sec:sha-typical-sha}

Usually, the SHA accelerator will process all blocks of a message and produce a message digest before starting the next message digest.

However, \chipname{} SHA working in Typical SHA mode also supports optional ``interleaved" message digest calculation. Users can insert new calculation (both Typical SHA and DMA-SHA) each time the SHA accelerator completes one message block. To be more specific, users can store the message digest in registers \hyperref[regdesc:SHAHNREG]{SHA\_H\_\regindex{n}\_REG} after completing each message block, and assign the accelerator with other higher priority tasks. After the inserted calculation completes, users can put the message digest stored back to registers \hyperref[regdesc:SHAHNREG]{SHA\_H\_\regindex{n}\_REG}, and resume the accelerator with the previously paused calculation.

\subparagraph{Typical SHA Process (except for SHA-512/{\it{t}})}

    \subfile{\modulefiles/27-sha-typical-process-except-512t__EN}

\subparagraph{Typical SHA Process (SHA-512/{\it{t}})}

    \subfile{\modulefiles/27-sha-typical-process-512t__EN}

\begin{tiplisting}\label{others:sha-note-typical-sha}
\vspace{-2.5em}
\begin{enumerate}
\item In this step, the software can also write the next message block (to be processed) in registers \hyperref[regdesc:SHAMNREG]{SHA\_M\_\regindex{n}\_REG}, if any, while the hardware starts SHA calculation, to save time.
\item You are resuming the SHA accelerator with the previously paused calculation.
\item Here you can decide if you want to insert other calculations. If yes, please go to the \hyperref[others:sha-interleaved]{process for interleaved calculations} for details.
\end{enumerate}
\end{tiplisting}

\subfile{\modulefiles/27-SHA-typical-process-interleaved__EN}


\subsubsection{DMA-SHA Mode Process}\label{sec:sha-dma-sha}

\chipname{} SHA accelerator does not support ``interleaving" message digest calculation when using the DMA, which means you cannot insert new calculation before the whole DMA-SHA process completes. In this case, users who need inserted calculation are recommended to divide your message blocks and perform several DMA-SHA calculation, instead of trying to compute all the messages in one go.

In contrast to the Typical SHA working mode, when the SHA accelerator is working under the DMA-SHA mode, all data read are completed via DMA.

Therefore, users are required to configure the DMA controller following the description in Chapter \ref{mod:dma} \textit{\nameref{mod:dma}}.

\subparagraph{DMA-SHA process (except SHA-512/{\it{t}})}

    \subfile{\modulefiles/27-sha-dma-process-except-512t__EN}

\subparagraph{DMA-SHA process for SHA-512/{\it{t}}}

    \subfile{\modulefiles/27-sha-dma-process-512t__EN}



\subsection{Message Digest}

After the hash task completes, the SHA accelerator writes the message digest from the task to registers \hyperref[regdesc:SHAHNREG]{SHA\_H\_\regindex{n}\_REG}(\regindex{n}: 0\verb+~+15). The lengths of the generated message digest are different depending on different hash algorithms. For details, see Table \ref{tab:sha-message-digest} below:

\begin{table}[H]
    \centering
    \caption{The Storage and Length of Message digest from Different Algorithms}
    \label{tab:sha-message-digest}

    \begin{threeparttable}
    \begin{tabular}{| l | c | l | }
    \hline
        \rowcolor{lightgray}
        Hash Algorithm & Length of Message Digest (in bits) & Storage\tnote{1} \\\hline

        % table contents
        SHA-1 & 160 & SHA\_H\_0\_REG\verb+ ~ +SHA\_H\_4\_REG\\
        \hline
        SHA-224 & 224 & SHA\_H\_0\_REG\verb+ ~ +SHA\_H\_6\_REG\\
        \hline
        SHA-256 & 256 & SHA\_H\_0\_REG\verb+ ~ +SHA\_H\_7\_REG\\
        \hline
        SHA-384 & 384 & SHA\_H\_0\_REG\verb+ ~ +SHA\_H\_11\_REG\\
        \hline
        SHA-512 & 512 & SHA\_H\_0\_REG\verb+ ~ +SHA\_H\_15\_REG\\
        \hline
        SHA-512/224 & 224 & SHA\_H\_0\_REG\verb+ ~ +SHA\_H\_6\_REG\\
        \hline
        SHA-512/256 & 256 & SHA\_H\_0\_REG\verb+ ~ +SHA\_H\_7\_REG\\
        \hline
        SHA-512/\regindex{t}\tnote{2} & \regindex{t} & SHA\_H\_0\_REG\verb+ ~ +SHA\_H\_\regindex{x}\_REG\\
        \hline
    \end{tabular}

        \begin{tablenotes}
            \item[1] The message digest are stored in registers from most significant bits to the least significant bits, with the first word stored in register \hyperref[regdesc:SHAHNREG]{SHA\_H\_0\_REG} and the second word stored in register \hyperref[regdesc:SHAHNREG]{SHA\_H\_1\_REG}... For details, please see subsection \ref{sec:sha-parsing-msg}.
            \item[2] The registers used for SHA-512/{\it{t}} algorithm depend on the value of \begin{math}t\end{math}. {\it{x}}+1 indicates the number of 32-bit registers used to store {\it{t}} bits of message digest, so that {\it{x}} = \begin{math}roundup(t/32)\end{math}-1. For example:
                \begin{itemize}
                    \item When {\it{t}} = 8, then {\it{x}} = 0, indicating that the 8-bit long message digest is stored in the most significant 8 bits of register \hyperref[regdesc:SHAHNREG]{SHA\_H\_0\_REG};
                    \item When {\it{t}} = 32, then {\it{x}} = 0, indicating that the 32-bit long message digest is stored in register \hyperref[regdesc:SHAHNREG]{SHA\_H\_0\_REG};
                    \item When {\it{t}} = 132, then {\it{x}} = 4, indicating that the 132-bit long message digest is stored in registers \hyperref[regdesc:SHAHNREG]{SHA\_H\_0\_REG}, \hyperref[regdesc:SHAHNREG]{SHA\_H\_1\_REG}, \hyperref[regdesc:SHAHNREG]{SHA\_H\_2\_REG}, \hyperref[regdesc:SHAHNREG]{SHA\_H\_3\_REG}, and \hyperref[regdesc:SHAHNREG]{SHA\_H\_4\_REG}.
                \end{itemize}
        \end{tablenotes}
     \end{threeparttable}
    \end{table}

\clearpage
\subsection{Interrupt}

SHA accelerator supports interrupt on the completion of message digest calculation when working in the DMA-SHA mode. To enable this function, write 1 to register \hyperref[regdesc:SHAINTENAREG]{SHA\_INT\_ENA\_REG}. Note that the interrupt should be cleared by software after use via setting the \hyperref[regdesc:SHAINTCLEARREG]{SHA\_INT\_CLEAR\_REG} register to 1.

\section{Register Summary} \label{sec:sha-reg-summ} \hypertarget{sha-reg-summ}{}

The addresses in this section are relative to the SHA accelerator base address provided in Table \ref{tab:sysmem-base-address} in Chapter \ref{mod:sysmem} \textit{\nameref{mod:sysmem}}.

The abbreviations given in Column \textbf{Access} are explained in Section \hyperref[glossary-access-types]{\textit{Access Types for Registers}}.

\subfile{\modulefiles/27-SHA-reg-summ__EN}

\section{Registers}

The addresses in this section are relative to the SHA accelerator base address provided in Table \ref{tab:sysmem-base-address} in Chapter \ref{mod:sysmem} \textit{\nameref{mod:sysmem}}.

\subfile{\modulefiles/27-SHA-reg__EN}


\end{document}
